\section{Research motivation}
\label{sec:motivation}

\subsection{The accessibility gap}

Security practitioners increasingly encounter AI assistants that assume
connectivity to a hosted provider. This creates at least four barriers. First,
\emph{cost}: training and inference are priced for organizations, not
individuals. Second, \emph{hardware}: full fine-tuning of even a few billion
parameters can exceed consumer VRAM. Third, \emph{privacy}: security telemetry,
source code, and incident data are sensitive and may be regulated. Fourth,
\emph{availability}: offline and air-gapped environments cannot depend on an
external API.

A sub-billion-parameter model will not match the broad capability of a
frontier model. The relevant question is whether specialization, retrieval, and
careful post-training can make a small model useful for bounded defensive tasks
at a fraction of the cost. If the answer is yes for even a subset of tasks, the
result is valuable for education, local tooling, and reproducible research. If
the answer is no, the negative result is also informative: it would quantify the
capability floor of current sub-billion models under realistic constraints.

\subsection{Why cybersecurity?}

Cybersecurity is a challenging but tractable specialization domain. It has
authoritative, structured public knowledge (ATT\&CK, CWE, NVD, OWASP, vendor
advisories), strong conventions for code and command syntax, and tasks with
verifiable answers such as vulnerability classification or secure-code
completion. It also has high safety stakes: a fluent but wrong answer can lead
to a misconfiguration or an unsafe command. This makes honest evaluation and
non-execution defaults especially important.

\subsection{Design priorities}

The project optimizes for the following, in order:

\begin{enumerate}[leftmargin=1.4em]
  \item \textbf{Low training cost.} A single consumer GPU or CPU should be
  enough to run the pipeline at small scale.
  \item \textbf{Low VRAM.} QLoRA, gradient checkpointing, short sequences, and
  measured memory reporting are first-class.
  \item \textbf{Fast, local inference.} CPU, GPU, and MPS paths; streaming;
  optional quantization; no provider dependency.
  \item \textbf{Offline operation.} After model download, inference and
  evaluation require no network.
  \item \textbf{Domain specialization.} Cybersecurity data, secure coding,
  detection engineering, and defensive workflows.
  \item \textbf{Context-aware autocomplete.} Bash, Python, PowerShell, YAML
  rules, regular expressions, filters, and configuration files.
  \item \textbf{Reproducibility.} Seeds, hashes, manifests, pinned environment
  records, and raw predictions.
  \item \textbf{Accessibility.} Apache-2.0 code, explicit third-party licenses,
  and a workflow that does not require paid services.
\end{enumerate}

\subsection{Target hardware as an experimental variable}

The project targets 4, 6, 8, 12, and 16 GB VRAM classes, plus CPU-only
deployment. These are experimental targets rather than guaranteed requirements.
Memory use depends on the attention kernel, sequence length, batch size,
optimizer, allocator, and framework version. The pipeline therefore records
measured peak memory next to each run instead of asserting that a configuration
``fits.''
